\chapter{CLI 与 Daemon 完整命令参考}
\label{ch:daemon-cli-v4}

\section{概述}

ebbackup 用户操作面由三部分组成：
\begin{itemize}[nosep]
  \item \textbf{CLI} \texttt{eb}：\filepath{engine\_cpp/cli/eb\_main.cc}，命令行入口；
  \item \textbf{Daemon 调度}：\texttt{schedule} / \texttt{watch}，\filepath{src/daemon/backup\_daemon.cc}；
  \item \textbf{C API}：\filepath{src/engine/eb\_backup\_capi.cc}，稳定 ABI（本章聚焦 CLI/daemon）。
\end{itemize}

除非注明，路径参数均为本地文件系统路径；加密仓库通过 \texttt{--password-env} 或 \texttt{--password-file} 提供口令。

\section{eb 命令总表}

\begin{longtable}{@{}p{2.2cm}p{4.8cm}p{7.2cm}@{}}
\toprule
命令 & 标志 / 参数 & 说明 \\
\midrule
\endfirsthead
\toprule
命令 & 标志 / 参数 & 说明 \\
\midrule
\endhead

\texttt{init} &
\texttt{<repo>} &
创建新仓库目录结构 \\
&
\texttt{--legacy-init} &
使用 \texttt{InitV03Repo}：legacy chunks，无 EbPack/snapshot 等 \\
&
\texttt{--legacy-digest} &
Legacy SHA256 而非 Standard digest \\
&
\texttt{--encrypt} &
（若实现）初始化加密 salt \\
\midrule

\texttt{backup} &
\texttt{<repo> <source>} &
运行 full 或 incremental backup \\
&
\texttt{--incremental} &
增量模式（需 prior manifest） \\
&
\texttt{--progress} &
stderr 进度 permille（0--1000） \\
&
\texttt{--lz4} &
启用 LZ4 压缩 \\
&
\texttt{--compress auto|lz4|zstd|off} &
压缩模式；\texttt{auto} 默认 cpu-budget=600‰ \\
&
\texttt{--cpu-budget PCT} &
1--100，映射 \texttt{cpu\_budget\_permille=PCT*10} \\
&
\texttt{--durability strict|balanced} &
EbPack flush/fsync 策略 \\
&
\texttt{--pipeline} / \texttt{--no-pipeline} &
强制启用 / 禁用 pipeline \\
&
\texttt{--encrypt} &
内容 AES-256-GCM 加密 \\
&
\texttt{--password-env VAR} / \texttt{--password-file PATH} &
加密口令来源 \\
&
\texttt{--filter-file PATH} &
filter 键值文件 \\
&
\texttt{--include PATH} / \texttt{--exclude PATH} &
路径前缀 filter \\
&
\texttt{--include-glob GLOB} / \texttt{--exclude-glob GLOB} &
glob filter \\
&
\texttt{--ext EXT} &
扩展名 filter \\
&
\texttt{--min-size N} / \texttt{--max-size N} &
文件大小（字节） \\
&
\texttt{--mtime-after UNIX} / \texttt{--mtime-before UNIX} &
mtime 窗口 \\
&
\texttt{--uid N} &
uid 精确匹配 \\
\midrule

\texttt{verify} &
\texttt{<repo>} &
仓库一致性校验 \\
&
\texttt{--at TXN} / \texttt{--at latest} &
VerifyAt snapshot（\texttt{latest}=0 当前） \\
&
\texttt{--require-anchor} &
强制 CARL anchor 验证 \\
&
\texttt{--password-env} / \texttt{--password-file} &
加密仓库解密 \\
\midrule

\texttt{recover} &
\texttt{<repo>} &
修复中断 txn $\rightarrow$ \texttt{kAborted} \\
\midrule

\texttt{restore} &
\texttt{<repo> <dest>} &
从 manifest 恢复文件树 \\
&
\texttt{--at TXN} &
时间旅行 snapshot \\
&
\texttt{--filter-file}、include/exclude/glob/ext/size/mtime/uid &
与 backup 相同 filter 语义 \\
&
\texttt{--verify-content} &
恢复后流式 chunk 内容校验 \\
&
\texttt{--skip-content-verify} &
跳过子集 Merkle 与内容校验 \\
&
\texttt{--password-env} / \texttt{--password-file} &
加密仓库 \\
\midrule

\texttt{export} &
\texttt{<repo> <bundle.ebb>} &
导出 bundle \\
&
\texttt{--encrypt-bundle} &
bundle 加密 \\
\midrule

\texttt{import} &
\texttt{<bundle.ebb> <repo>} &
导入 bundle 至新仓库 \\
\midrule

\texttt{schedule} &
\texttt{<config>} &
Daemon 定时 incremental \\
&
\texttt{--once} &
仅运行一轮 \\
\midrule

\texttt{watch} &
\texttt{<repo> <source>} &
目录变更触发 incremental \\
&
\texttt{--debounce-ms N} &
去抖毫秒，默认 2000 \\
&
\texttt{--once} &
仅触发一次 \\
&
backup 同名 flags &
lz4/pipeline/encrypt/filter 等 \\
\midrule

\texttt{gc-orphans} &
\texttt{<repo>} &
回收无引用 chunk \\
&
\texttt{--dry-run} &
仅统计 \\
&
\texttt{--latest-only} &
仅 latest manifest 引用的 hash（默认含全部保留 snapshot） \\
\midrule

\texttt{compact} &
\texttt{<repo>} &
EbPack 重写压缩 \\
&
\texttt{--dry-run} &
预览 \\
&
\texttt{--wait-idle SEC} &
等待仓库 idle 再 compact \\
\midrule

\texttt{repo-stats} &
\texttt{<repo>} &
物理/live/orphan/ampl\_ratio 等指标 \\
\midrule

\texttt{list-snapshots} &
\texttt{<repo>} &
列出 snapshot \\
&
\texttt{--json} &
JSON 数组输出 \\
\midrule

\texttt{prune-snapshots} &
\texttt{<repo>} &
GFS 保留策略 prune \\
&
\texttt{--dry-run} &
预览 \\
&
\texttt{--retain-min N} &
最少保留 snapshot 数 \\
&
\texttt{--retention-tiers SPEC} &
如 \texttt{1h:24,1d:7,7d:4,30d:6} \\
\midrule

\texttt{bench cdc} &
\texttt{<file>} &
FastCdcSlice 吞吐 micro-bench \\
\midrule

\texttt{bench hcrbo} &
\texttt{<file> <delta\_offset>} &
HCRBO 增量 reuse micro-bench \\
\bottomrule
\caption{eb CLI 完整命令与标志}
\label{tab:eb-cli-full}
\end{longtable}

退出码：成功 0；失败打印 \texttt{error: <message>} 返回 1。backup 成功额外打印 \texttt{stats: files=... chunks\_written=...}。

\section{Init 默认 Feature Bitmask}

\texttt{eb init <repo>}（无 \texttt{--legacy-init}）调用 \texttt{BackupEngine::InitRepoEx}，等价 C API \texttt{eb\_backup\_init\_repo\_ex(0)}，默认启用：

\begin{table}[htbp]
\centering
\caption{Init 默认 feature 组合（v0.6+）}
\begin{tabular}{@{}ll@{}}
\toprule
Feature & 含义 \\
\midrule
\texttt{persistent\_index} & \filepath{chunk.idx} 持久化索引 \\
\texttt{manifest\_binary}（manifest v4） & 32B hash 二进制 manifest \\
\texttt{compress-auto} & 内容分类自适应压缩 \\
\texttt{snapshots} & \filepath{snapshots/} 时间旅行 \\
\texttt{EbPack} & 16-shard pack 存储 \\
\texttt{coalesced\_meta} & backup 进行中延迟 superblock 落盘 \\
\bottomrule
\end{tabular}
\end{table}

\texttt{--legacy-init} 回退 \texttt{InitV03Repo()}：legacy \filepath{data/chunks/} append，无上述现代 feature。

\texttt{--legacy-digest} 设置 \texttt{standard\_digest=false}，使用 Legacy SHA256 内容 hash。

\section{Schedule JSON 与键值配置}

\subsection{ScheduleConfig 结构}

\begin{structbox}[ScheduleConfig 字段]
\begin{itemize}[nosep]
  \item \texttt{interval\_seconds}：默认 3600，轮询间隔；
  \item \texttt{source\_path}：备份源目录（\textbf{必填}）；
  \item \texttt{repo\_base}：仓库基目录（\textbf{必填}），实际 repo = \texttt{repo\_base/current}；
  \item \texttt{retain\_count}：默认 3，legacy \texttt{repo-*} 旋转保留数；
  \item \texttt{retention\_policy}：GFS tiers + \texttt{retain\_min}；
  \item \texttt{auto\_prune}：默认 true，每轮 backup 后 prune snapshots；
  \item \texttt{auto\_gc\_after\_prune}：默认 false，prune 后 \texttt{GcOrphans}；
  \item \texttt{backup\_options}：嵌套 \texttt{BackupOptions}（compress/pipeline/encrypt/filter 等）；
  \item \texttt{encryption\_password}：schedule 级口令（亦写入 backup\_options）。
\end{itemize}
\end{structbox}

\subsection{JSON Schema 字段}

\texttt{LoadScheduleConfigAuto}：若文件首非空白字符为 \texttt{\{}，按 JSON 解析；否则按 \texttt{key=value} 行解析。JSON 支持的键（与 INI 键同名）：

\begin{table}[htbp]
\centering
\caption{Schedule 配置键}
\small
\begin{tabular}{@{}llp{6cm}@{}}
\toprule
键 & 类型 & 说明 \\
\midrule
\texttt{interval\_seconds} & int &  sleep 间隔 \\
\texttt{source} & string & 源路径 \\
\texttt{repo\_base} & string & 仓库基目录 \\
\texttt{retain} / \texttt{retain\_min} & int & 保留 snapshot 下限 \\
\texttt{retention\_tiers} & string & GFS tier 规格 \\
\texttt{auto\_prune} & bool & \texttt{true}/\texttt{1}/\texttt{yes} \\
\texttt{auto\_gc\_after\_prune} & bool & prune 后 GC \\
\texttt{lz4} / \texttt{pipeline} / \texttt{encrypt} & bool & backup 选项 \\
\texttt{password} & string & 加密口令 \\
\texttt{filter\_file} & path & 加载 filter 文件 \\
\texttt{include\_glob} / \texttt{exclude\_glob} & string & 追加 glob \\
\texttt{compress} & string & \texttt{auto|lz4|zstd|off} \\
\texttt{cpu\_budget} & int & 1--100 百分比 \\
\texttt{durability} & string & \texttt{strict|balanced} \\
\bottomrule
\end{tabular}
\end{table}

\subsection{JSON 示例}

\begin{lstlisting}[caption={schedule.json 示例},label={lst:schedule-json}]
{
  "interval_seconds": 1800,
  "source": "/data/projects",
  "repo_base": "/backup/eb",
  "retain_min": 5,
  "retention_tiers": "1h:24,1d:7,7d:4,30d:6",
  "auto_prune": true,
  "auto_gc_after_prune": true,
  "compress": "auto",
  "cpu_budget": 60,
  "durability": "balanced",
  "pipeline": true,
  "encrypt": false,
  "filter_file": "/etc/ebbackup/filter.conf",
  "exclude_glob": "*.tmp"
}
\end{lstlisting}

\subsection{INI 示例}

\begin{lstlisting}[caption={schedule.conf 示例},label={lst:schedule-ini}]
interval_seconds=3600
source=/home/user/Documents
repo_base=/mnt/backup/eb
retain_min=3
auto_prune=true
compress=auto
cpu_budget=80
pipeline=true
filter_file=/home/user/.ebbackup/filter
\end{lstlisting}

\subsection{RunScheduledBackup 行为}

\begin{enumerate}
  \item \texttt{repo = repo\_base/current}，不存在则 \texttt{InitRepoEx}（现代 feature）；
  \item 循环：首轮 full，后续 incremental；
  \item 可选 \texttt{PruneSnapshots} + \texttt{GcOrphans}；
  \item legacy \texttt{repo-YYYYMMDD-HHMMSS} 目录若存在则 \texttt{PruneRotatedRepos}；
  \item \texttt{--once} 或 \texttt{max\_runs=1} 仅一轮；否则 sleep \texttt{interval\_seconds}。
\end{enumerate}

\section{Watch 模式}

\texttt{eb watch <repo> <source>} 调用 \texttt{RunWatchBackup}：

\begin{itemize}[nosep]
  \item \texttt{FsWatch}：Linux inotify / Windows ReadDirectoryChangesW；
  \item \texttt{WaitForChange(debounce\_ms)}：默认 2000ms 去抖，合并 burst 事件；
  \item 触发 \texttt{RunOneBackup(repo, source, kIncremental, options)}；
  \item \texttt{--once}：单次触发后退出；否则无限循环。
\end{itemize}

Watch 不自动 prune/GC（与 schedule 不同），适合开发目录实时保护。repo 需已 init；首次 backup 建议手动 full。

\section{Filter 文件格式}

\subsection{语法}

\begin{itemize}[nosep]
  \item 每行 \texttt{key=value}，首尾空白 trim；
  \item \texttt{\#} 开头为注释；
  \item 同一 key 可重复（如多条 \texttt{exclude\_glob}）；
  \item 键名见 \texttt{LoadBackupFilterFromFile}（\filepath{backup\_filter.cc}）。
\end{itemize}

\subsection{支持的键}

\begin{table}[htbp]
\centering
\caption{Filter 文件键}
\begin{tabular}{@{}ll@{}}
\toprule
键 & 映射字段 \\
\midrule
\texttt{include\_path} & \texttt{include\_paths[]} \\
\texttt{exclude\_path} & \texttt{exclude\_paths[]} \\
\texttt{include\_glob} & \texttt{include\_globs[]} \\
\texttt{exclude\_glob} / \texttt{name\_glob} & \texttt{exclude\_globs[]} \\
\texttt{ext} / \texttt{extension} & \texttt{extensions[]} \\
\texttt{min\_size} / \texttt{max\_size} & uint64 \\
\texttt{mtime\_after} / \texttt{mtime\_before} & int64 Unix \\
\texttt{uid} & uint32 \\
\bottomrule
\end{tabular}
\end{table}

\subsection{完整示例}

\begin{lstlisting}[caption={生产环境 filter 示例},label={lst:filter-prod}]
# === ebbackup filter ===
# 排除缓存与依赖
exclude_path=node_modules
exclude_path=__pycache__
exclude_path=.cache
exclude_glob=*.pyc
exclude_glob=*.swp
exclude_glob=*.tmp
exclude_glob=.DS_Store

# 日志仅保留最近（配合 mtime）
exclude_glob=*.log
mtime_before=1893456000

# 大型虚拟机磁盘（按扩展名）
exclude_glob=*.vmdk
exclude_glob=*.iso

# 可选：仅包含源码树
# include_path=src
# include_glob=*.rs
# include_glob=*.toml

min_size=1
max_size=5368709120
\end{lstlisting}

CLI \texttt{--filter-file} 与 schedule \texttt{filter\_file}、watch 均调用同一 loader；C API \texttt{eb\_backup\_load\_filter\_file} 语法一致。

\section{密码与加密 CLI 惯例}

\texttt{ReadPassword} 优先级：
\begin{enumerate}
  \item \texttt{--password-env VAR}：读取 \texttt{getenv(VAR)}；
  \item \texttt{--password-file PATH}：读首行；
  \item 否则空（非加密仓库 OK，加密仓库 backup/restore/verify 失败）。
\end{enumerate}

\section{Bench 子命令}

\begin{itemize}[nosep]
  \item \texttt{eb bench cdc FILE}：整文件读入内存，\texttt{FastCdcSlice::Chunk} 计时，输出 MB/s 与 avg chunk size；
  \item \texttt{eb bench hcrbo FILE OFFSET}：在 OFFSET 异或 0x5A 后跑 HCRBO incremental，输出 reuse\% 与 rolling skip hits。
\end{itemize}

\section{与环境变量的关系}

CLI 未直接暴露但影响引擎行为的环境变量（pipeline/CDC 见专门章节）：

\begin{table}[htbp]
\centering
\caption{CLI 相关的环境变量}
\small
\begin{tabular}{@{}ll@{}}
\toprule
变量 & 影响 \\
\midrule
\texttt{EBBACKUP\_PIPELINE\_WORKERS} & backup pipeline 并发 \\
\texttt{EBBACKUP\_PIPELINE\_PROFILE} & backup 结束 phase 打印 \\
\texttt{EBBACKUP\_CDC\_FAST\_SLICE} & FastCdcSlice 路由 \\
\texttt{EBBACKUP\_AUDIT\_KEY} & verify 时 CARL 签名验证 \\
\bottomrule
\end{tabular}
\end{table}

\section{典型工作流}

\subsection{首次部署}

\begin{lstlisting}[caption={初始化与首次 full backup},label={lst:workflow-init}]
eb init /backup/repo
eb backup /backup/repo /data/source --compress auto --progress
eb verify /backup/repo
\end{lstlisting}

\subsection{定时增量}

\begin{lstlisting}[caption={systemd 调用 schedule},label={lst:workflow-schedule}]
eb schedule /etc/ebbackup/schedule.json
\end{lstlisting}

\subsection{常驻服务（Wave O+）}

\textbf{Windows}：以管理员运行 \texttt{eb service install --config PATH} 注册 SCM 服务；\texttt{eb service run} 作为 ServiceMain 入口执行 \texttt{RunScheduleConfig}；\texttt{scripts/install\_service.ps1} 封装安装。Daemon 循环支持 \texttt{RequestDaemonStop()} 分片 sleep 优雅退出。

\textbf{Linux systemd}：模板 \filepath{deploy/ebbackup.service} + \filepath{deploy/ebbackup.timer}；\texttt{ExecStart=eb schedule /etc/ebbackup/queue\_drain.json --once}；配合 timer 周期性 drain job 队列。

\begin{lstlisting}[caption={eb service 子命令},label={lst:eb-service}]
eb service run --config /etc/ebbackup/queue_drain.json
eb service install --config /etc/ebbackup/queue_drain.json --name EbbackupDaemon
eb service status --name EbbackupDaemon
eb service uninstall --name EbbackupDaemon
\end{lstlisting}

\subsection{垂直插件（Wave P / ABI v27）}

\begin{lstlisting}[caption={插件列表与备份},label={lst:eb-plugins}]
eb plugin list
eb backup /backup/repo /data/source --plugins sqlite_checkpoint
eb job add /backup/repo --json '{"id":"db","name":"DB","source_path":"/data","plugins":["sqlite_checkpoint"]}'
\end{lstlisting}

Schedule / daemon JSON 或 INI 同样支持逗号分隔 \texttt{plugins}（与 \texttt{job\_ids} 风格一致）：

\begin{lstlisting}[caption={schedule 插件配置},label={lst:schedule-plugins}]
# schedule.ini
plugins=sqlite_checkpoint,registry_hive

# queue_drain.json
{"mode":"queue_drain","repo_path":"/backup/repo","plugins":"sqlite_checkpoint"}
\end{lstlisting}

备份完成后 \filepath{reports/<txn>.json} 含可选 \texttt{"plugins":[...]} 各插件运行摘要。

\subsection{智能排除建议（Wave R / ABI v28）}

\begin{lstlisting}[caption={排除建议与作业},label={lst:eb-suggest-excludes}]
eb suggest-excludes /data/source --json
eb suggest-excludes /data/source --filter-file /etc/ebbackup/filter.conf --include-ide
eb job add /backup/repo --json '{"id":"dev","name":"Dev","source_path":"/data/source","exclude_paths":["/data/source/node_modules"],"exclude_globs":["Thumbs.db"]}'
\end{lstlisting}

C API \texttt{eb\_backup\_suggest\_exclude\_filters\_json} 接受 \texttt{source\_path}、可选 \texttt{max\_depth}、\texttt{include\_ide\_dirs} 与 \texttt{existing} filter；返回 \texttt{items[]} 与 \texttt{recommended}。建议须用户确认后写入 \texttt{exclude\_paths} / \texttt{exclude\_globs} 或会话 filter。

\subsection{备份窗口（Wave T / ABI v29）}

\begin{lstlisting}[caption={作业备份窗口},label={lst:eb-job-window}]
eb job add /backup/repo --json '{"id":"nightly","name":"Nightly","source_path":"/data/source","window_start":"02:00","window_end":"06:00","deadline_grace_seconds":300,"durability_adaptive":true}'
\end{lstlisting}

窗外 \texttt{RunJob} 返回 \texttt{outside backup window}；队列 drain 跳过。距 \texttt{window\_end} 不足 grace 秒时 durability 降级为 Balanced；超时截断写入报告 \texttt{window\_truncated}。

\subsection{选择性恢复}

\begin{lstlisting}[caption={filter 恢复 + 内容校验},label={lst:workflow-restore}]
eb restore /backup/repo /restore/target \
  --filter-file /etc/ebbackup/filter.conf \
  --verify-content --at 128
\end{lstlisting}

\section{本章小结}

\texttt{eb} CLI 覆盖仓库生命周期：init、backup/restore、verify/recover、snapshot 运维、compact/GC、bundle 导入导出与 micro-bench。Init 默认启用 persistent index、manifest v4、snapshots、EbPack 与 coalesced meta。Daemon \texttt{schedule} 支持 JSON/INI 双格式，单 repo \filepath{current} 增量链；\texttt{watch} 提供 debounce 目录监视。Filter 文件与 CLI flags 共用 \texttt{BackupFilterOptions} 语义，backup 与 restore 一致。
